\section{Conclusion}
NetSuRF connects exact INT8 kernel-bit faults to budgeted logical protection and
task-level evaluation through a reproducible software flow. On the three tested
models, ranking quality is not a single property: association with isolated
damage, protection benefit, and computation cost can disagree. Per-example
absolute gradients improve several ResNet-8 allocations; KWS often admits a
competitive MSB baseline; ToyCar demonstrates why reconstruction-MSE sensitivity
cannot substitute for file-level AUC evidence. FKeras's strict bit-tier ordering
also shows why different scores need not produce different masks.

An exact counterexample establishes that isolated susceptibility has no universal
optimal-protection guarantee under simultaneous faults. The independent-pool
measurements do not, however, resolve a practical advantage of bounded
conditional selection over exact isolated allocation. Together, these positive,
negative and unresolved results support NetSuRF as a software methodology for
testing the reliability of ranking-driven decisions, not a universally superior
score or optimizer. Stronger partial-tier baselines, additional fault fields and
independent task populations are the next empirical steps; no physical mitigation
performance is claimed.
